\section{Related Work}
\label{sec:related}

\paragraph{Multilingual STR and Indic benchmarks.}
STR is mature for Latin script~\cite{parseq}, but coverage collapses elsewhere.
IndicSTR12~\cite{indicstr12} and the Bharat Scene Text Dataset~\cite{bstd} established real-image
benchmarks for Indian languages; on the latter, PARSeq reaches $\parseqSynth\%$ average WRR
trained on synthetic data only and $\parseqFinetuned\%$ after fine-tuning on real labeled images
--- numbers that presuppose labeled target data. GlotOCR~\cite{glotocr} quantifies the wider
problem: frontier VLMs read fewer than thirty of 100+ scripts. Synthetic-data engines for STR
remain Latin-centric~\cite{unionst}, and fine-tuning VLMs on synthetic renderings has been shown
for a single low-resource script with real-data validation~\cite{manchu}. Our work targets the
regime where no real labeled target image exists, and asks not only whether such scripts can be
read but how well, predicted in advance.

\paragraph{Tokenization granularity for complex scripts.}
MGP-STR~\cite{mgpstr} fuses character, BPE, and WordPiece streams within one Latin recognizer;
grapheme-level tokenization beats byte-BPE for Bengali handwriting~\cite{gradet}. In the LLM
literature, fertility's value as a tokenizer metric is contested~\cite{beyondfertility}. Our
preliminary study (Sec.~\ref{sec:method}) makes the supervised grapheme-vs-BPE choice
\emph{predictable} across nine scripts (polarity \lawPolarity, leave-one-script-out
$\lawLoso\%$), and this paper then tests --- and refutes --- the natural conjecture that the same
property predicts transfer.

\paragraph{Unseen units vs.\ unseen scripts.}
A substantial line reads unseen \emph{characters within a known script} by decomposition into
shared sub-units: stroke- and radical-based zero-shot Chinese
recognition~\cite{strokezs,star,hiercode}, and open-set recognition of novel
characters~\cite{openset1,openset2,moose}. Zero-shot \emph{detection} of unseen scripts localizes
but does not read~\cite{unseendet}. Byte-level pivots let ASR cover languages with unseen scripts
in speech~\cite{maestrou}, and task-arithmetic analogies achieve zero-real-data synthetic-to-real
HTR --- for five \emph{Latin-script} languages, with unseen writing systems explicitly left
open~\cite{taskanalogies}. To our knowledge, no prior system reads real scene text of an entirely
unseen script --- different Unicode block, different glyphs --- from zero real target images, let
alone with accuracy forecast beforehand.

\paragraph{What predicts cross-lingual transfer.}
In NLP, transfer tracks lexical and subword overlap rather than language
relatedness~\cite{subwordoverlap,overlap2}. For STR specifically, Ref.~\cite{crosslingstr}
concludes that source dataset size and visual similarity, not linguistic typology, drive
cross-lingual transfer. Our findings sharpen rather than contradict this, in a regime that work
did not study (zero real target images): (i) a typology-adjacent structural property ---
fertility --- indeed fails as a transfer predictor ($\spearFert$), as a data-centric view
expects; (ii) \emph{synthetic target exposure} dominates --- a target-side data-size effect;
(iii) within a fixed budget, coverage of the shared grapheme space orders transfer
(\spearCovEqual\ among equal-budget scripts) --- an OCR instantiation of the subword-overlap
principle in a \emph{designed} output space; and (iv) a frozen-encoder visual-similarity control
predicts nothing in our data (\spearVissim). Because LOSO balances source data by construction,
the confound identified by~\cite{crosslingstr} is controlled by design.

\paragraph{Prediction and preregistration in empirical ML.}
Scaling laws make supervised STR predictable in model and data size~\cite{strscaling};
preregistration and negative-result reporting are the reproducibility literature's recommended
remedies for HARKing and cherry-picking~\cite{preregml,repro}, yet full-study adoption is rare.
We run the entire campaign under a frozen preregistration with an append-only amendment log and
commit quantitative forecasts before the runs they predict; we offer the protocol itself as a
contribution.
